\subsection{目的、证明边界与失败语义}

本附录使审查者只依 PDF 中的数学定义与内嵌源码，即可独立重写有限
检查器。必须区分两类接口：定理~\ref{n5:thm:one-intersection-flag}
的模三端点穷尽是五阶定理的逻辑前提；其余有限集合、递推、Petersen
纤维轮廓、父集极值及小图割在正文中给出数学证明，相应程序只检查
抄录。旧带标号森林边只作交叉诊断。

必须区分：
\begin{enumerate}[label=(\arabic*)]
\item \emph{数学证书检查}：直接核查递推、纤维十项和、父集极值或图割；
\item \emph{独立实现复算}：不导入生产器，从定义重建同一有限对象；
\item \emph{文件身份检查}：SHA--256 与本附录清单一致。
\end{enumerate}
第三项不能代替第一项。任何缺失输入、未知字段、解析失败、情形数
不符、哈希不符或状态非 PASS，都必须使参考入口非零退出；不得读取
旧输出后直接打印 PASS。
PDF 内嵌的 \path{attachment_manifest.json} 给出其余每个附件的字节数
与 SHA--256；清单由构建脚本在嵌入前从实际字节生成。

\subsection{规范坐标与精确算术}

变量按 \(v(i,j)=5i+j\) 编号。三元组和二元组均按字典序排列：
\[
\begin{aligned}
\binom{[5]}3&=(012,013,014,023,024,034,123,124,134,234),\\
\binom{[5]}2&=(01,02,03,04,12,13,14,23,24,34).
\end{aligned}
\]
三次坐标点 \((R,C)\) 编号为 \(10\,\operatorname{ind}(R)+
\operatorname{ind}(C)\)，二次影子点同理。其九个可能父点是
\[
(A\cup\{r\},B\cup\{c\}),\qquad r\notin A,\ c\notin B.
\]
二次 permanent 基为
\[
q_{ij;ab}=x_{ia}x_{jb}+x_{ib}x_{ja},\qquad i<j,\ a<b.
\]

全部 JSON 采用 UTF--8（无 BOM）、LF 换行和文件末尾恰一个 LF。
整数位集不得经过浮点转换。有理数秩使用
\texttt{fractions.Fraction} 或 fraction-free Bareiss 消元。除下一节明确
列为定理前提的一维交端点证书外，五阶诊断不需要有限域。模素数矩阵秩
只能作为同一整数矩阵的特征零秩下界；等价地，模素数核维只能作为
特征零核维上界，不能把有限域等号无条件提升为特征零等号。

\subsection{\texorpdfstring{\(n=3,4\)}{n=3,4} 的独立接口}

按正文直接构造 \(E_n\) 和 Koszul 整数矩阵，应得到
\[
\begin{array}{c|rr}
&n=3&n=4\\ \hline
\dim E_n&9&36\\
\dim E_n^{(1)}&1&16\\
\rank K(\perm_n)&80&560\\
\text{独立因子单项秩}&26&92.
\end{array}
\]
四阶图表证书还应重建 \(659\times659\) 方形子式、Schur 常数
\(-32768\)，并检查十五个归一化联合图均为 DAG。参考命令为
\begin{verbatim}
python perm35_exact_verification.py --n 3 --psi
python perm4_rank8_independent_audit.py
python perm4_rank8_verify_all.py
\end{verbatim}

\subsection{一维交旗标证书的独立重建}
\label{app:one-intersection-flag-spec}

本接口严格核查定理~\ref{n5:thm:one-intersection-flag} 的有限端点。
变量编号为 \(v(i,a)=5i+a\)。\(Q=\Sym^2V/E\) 的规范有序基由下列
\(225\) 个描述符组成：
\[
\begin{array}{c|c|c}
\text{描述符}&\text{范围}&\text{代表元}\\ \hline
S(i,a)&0\le i,a<5&x_{ia}^2\\
R(i;a,b)&a<b&x_{ia}x_{ib}\\
C(i,j;a)&i<j&x_{ia}x_{ja}\\
X(i,j;a,b)&i<j,\ a<b&x_{ia}x_{jb}.
\end{array}
\]
在最后一类中，关系
\(x_{ia}x_{jb}+x_{ib}x_{ja}=0\) 位于 \(Q\)，故另一条对角线映为
\(-X(i,j;a,b)\)。这给出 \(25+50+50+100=225\) 个互异行列环面权。

三次 divided-power 基按非降三元组
\((u,v,w)\)、\(0\le u\le v\le w<25\) 编号，共
\(\binom{27}{3}=2925\) 个。若 \(m\) 是这样的多重集，则极化映射为
\begin{equation}\label{app:eq:divided-polarization}
 \Delta(m)=\sum_{r\in\operatorname{supp}(m)}
 x_r\otimes[m\setminus\{r\}]_Q.
\end{equation}
这里按 divided-power 约定，每个不同的 \(r\) 只出现一次；商类的符号
由上表决定。因此所有矩阵元都在 \(\{0,\pm1\}\)。按三次行列重数分块
应得到恰 \(1225\) 个块。对选择的商权集合 \(W\)，删去
式~\eqref{app:eq:divided-polarization} 中商权属于 \(W\) 的输出行，再在
\(\mathbf F_3\) 上作精确高斯消元。记所得约化矩阵为
\(\overline\Phi_W\)；这里不需要把 divided-power 格在特征三解释成
通常的对称幂。各块的核维增量之和就是
\[
 \overline p_3(W)=
 \dim_{\mathbf F_3}\ker\overline\Phi_W-
 \dim_{\mathbf F_3}\ker\overline\Phi_0.
\]
必须先独立得到
\(\dim_{\mathbf F_3}\ker\overline\Phi_0=100\)。

端点集合按以下规则生成，不得读取冻结 JSON：
\begin{enumerate}[label=(\arabic*)]
\item 枚举全部 \(100=\binom52^2\) 个坐标矩形。四维端点取其唯一九权
商像，并加入 \(225-9\) 个可能的第十权，共 \(21{,}600\) 个旗标。
\item 五维端点只需两个对称代表：
\[
 (0,1,2,5,6)\quad\text{(attached)},\qquad
 (0,1,5,6,12)\quad\text{(external)}.
\]
两者的 \(\Sym^2\) 商像均有十四权。对每个代表枚举全部
\(\binom{14}{9}\) 个九权子集，再加入 \(225-9\) 个可能的第十权，
每类 \(432{,}432\) 个旗标。独立枚举所有含矩形的坐标五格集还必须
得到 attached 类 \(1200\) 个、external 类 \(900\) 个；每个五格集
只含一个矩形。
\end{enumerate}

预期直方图必须逐项相等：
\[
\begin{array}{c|r}
\multicolumn{2}{c}{d=4}\\ \hline
p&\#\\ \hline
20&12300\\21&5700\\22&3600
\end{array}
\]
五维的合并直方图为
\[
\begin{array}{c|rrrrrr}
p&10&11&12&13&14&15\\ \hline
\#&2238&10854&44407&88827&143962&166486
\end{array}
\]
\[
\begin{array}{c|rrrrr}
p&16&17&18&19&20\\ \hline
\#&149800&110299&67615&37787&19798
\end{array}
\]
\[
\begin{array}{c|rrrrrr}
p&21&22&23&24&25&26\\ \hline
\#&11214&7394&2768&944&451&20
\end{array}
\]
五维直方图是两类之和；其分类最大值必须分别为 \(26\) 与 \(22\)。
程序还必须检查局部非零真值表项总数 \(43825\)，从而防止空枚举或漏块。

规范命令为
\begin{verbatim}
python -O perm5_one_intersection_flag_standalone_exact.py
python -O perm5_one_intersection_flag_standalone_exact.py \
  --output replay.json
\end{verbatim}
程序默认 compare-only；只有显式给出 \path{--output} 才写文件。任何
计数、直方图、最大值或内部不变量不符都必须非零退出，且不得依赖
Python 的 \texttt{assert}。内嵌源码与规范输出的 SHA--256 分别为
\begin{verbatim}
7578DB64A6C501230A1155F88453D6A3B99C6B58DBFEB3F87E2A5B8BD69DDFDE
B373461AB31F760C7A3DC2F83BFC61EAD73822C925FA8E81C5ACB015BD9705E0
\end{verbatim}
模三计算在这里给出的是特征零核维的上界：同一整数矩阵满足
\(\rank_{\mathbf F_3}\le\rank_{\Q}\)，故
\(\dim\ker_{\Q}\le\dim\ker_{\mathbf F_3}\)。它不是把有限域等号
无条件提升为特征零等号。

\subsection{五阶状态宇宙与小图表}

在进入状态宇宙前，先核查拆项入口。基础展平给秩至少十；若最小分解
有 \(10\le r\le15\) 项，则保留六个原始项，并反复使用
\(T=\frac12T+\frac12T\) 把其余 \(r-6\) 项补成恰九项。因此 fixed-six
反证直接覆盖全部 \(r\le15\)，不得再把旧 lower--15 SAT/DRAT 层列为
前置依赖。

Petersen 图顶点为 \(\binom{[5]}2\)，不交时相邻。直接检查
\[
n(k)=\min_{|A|=k}|N(A)|
=(0,3,5,6,6,8,9,9,10,10,10).
\]
正文的十层递推给出表~\eqref{n5:eq:shadow-table}。fixed-six 状态只由
\[
s\in\{19,20,21,22\},\quad m_s\le t,\quad d\ge9,\quad t+d\le60
\]
产生，故计数为
\[
\binom82+3\binom52=58.
\]
八条路由的计数必须为
\[
38+1+1+9+2+1+1+5=58.
\]
这两个等式是闭式整数计数，不是对 58 个状态的逐态穷举。

小图必须从坐标定义重建，而不能读取预存真值表。对每个三次环面权，
以该权的三次单项式为顶点；对每个变量与二次商权 \(q\)，把偏导系数
行的两个非零项画成带符号的 \(q\)-边，把单个非零项画成 \(q\)-锚。
选择 \(q\) 就删除所有 \(q\)-关系。四边形、六边图和 \(K_{3,3}\)
中的符号可由顶点换号消去，故核维数是未锚连通分量数。

固定一个 crossing 权 \(x\)，设已有 \(N\le11\) 个方向。审查者应由
四边形、六边图与 \(K_{3,3}\) 的 inclusion-minimal cuts 直接重建
\[
\Delta_xp\le
\begin{cases}
N,&0\le N\le4,\\
\lceil3N/4\rceil,&5\le N\le11.
\end{cases}
\]
非角部分的规范参数为
\[
\begin{gathered}
u=h_2+h_1,\quad v=v_2+v_1,\quad
\alpha=h_2,\quad\gamma=v_2,\\
P=u+v,\quad A=P+\alpha+\gamma,\quad
z=\alpha\gamma,\quad o=uv,
\end{gathered}
\]
并且对重复块数 \(r\)、matching 块数 \(m\) 和非角预算 \(q\)，必须
逐割得到
\[
r\le2P,\qquad q\ge Q:=r+A+D(m;z,o),
\]
\[
D(m;z,o)=
\begin{cases}
0,&m\le z,\\
m-z,&z<m\le o,\\
2m-o-z,&m>o.
\end{cases}
\]
其中六个重复组各含两条 side crossing；每个组的两个六边块需要
互不相交的 external 标签。matching 格的六个顶点是三行三列的六个
完美匹配，九条 \(K_{3,3}\) 边以 \(2\times2\) 矩形标号。删去
\(x\) 后，横纵组都 double-side、都 side-positive（扣除前者）、其余三种情形
分别给 diagonal 下界 \(0,1,2\)，且不同格的 diagonal 标签不交。
在 \(Q\le11\) 下，三段分别推出 \(4(r+m)\le3Q\)，并另证
\(F(4)\le2\)。角块满足
\[
g(b)=\lfloor b^2/4\rfloor,\qquad
c\le\min\{4,s+g(b)\}.
\]
除 \((s,b)=(4,0),(0,4)\) 外有
\(c\le\lceil3(s+b)/4\rceil\)；两个例外由 \(F(4)\le2\) 闭合。
因此旧十二项 crossing 边际表不再是活跃证明依赖。

无 crossing 权集合不得再调用旧 square 轨道表。审查者应从正文的
六类模式计数重建行、列二点压缩；核心局部不等式是
\[
(x\vee y)(a\vee b)+(x\wedge y)(a\wedge b)\ge xa+yb.
\]
压缩后 \(G_0\supseteq\cdots\supseteq G_4\)、
\(H_0\supseteq\cdots\supseteq H_4\)，square 集为 Ferrers ideal。
五图三角形包络必须由 Kruskal--Katona 公式重建为
\[
\tau(e)=(0,0,0,1,1,2,4,4,5,7,10,10,10)_e.
\]
若 square、row edge、column edge 数为 (s,r,c)，则必须直接验证
\[
Q\le s+R_0+C_0+\min\{rc,4R_0,4C_0\},\qquad
R_0=\min(4s,2r),\quad C_0=\min(4s,2c).
\]
正文列出的十种粗界例外只能使用 one-square、row-domino、row-triple
和 L 形的闭式度数界关闭。由此无 crossing 层必须重现
\[
B_9\le35,\qquad B_{10}\le50,\qquad
B_{11}\le55,\qquad B_{12}\le60,
\]
且九维等号必须恰为一百个 pure row/column \(K_5-e\)。数学展开见
\path{n5_nocrossing_compression_structural_reduction_20260810.md}；
\path{perm5_nocrossing_compression_diagnostic.py} 的 shifted 状态遍历仅作
反例诊断，不进入证明依赖。与上述闭式 crossing 边际定理相加必须重现
\[
p_9\le35,\qquad p_{10}\le50,\qquad p_{11}\le55,\qquad p_{12}\le61.
\]
这里不得遗漏密度界的等号余量：一层 pure \(K_5\) 不含 square、
side crossing 或 diagonal crossing，只能提供同一方向的
boundary/external 标签，故加入任意 crossing 的边际为零；再加入一个
square 后，非角块仍不激活，
四个 corner 中至多激活一个，故边际至多一。十一维按 crossing 数
\(0,1,\ge2\) 分层；十二维按 \(0,1,2,\ge3\) 分层。最后一层
\(3,\ldots,12\) 个 crossing 的粗值必须为
\[
59,58,60,61,56,57,57,58,56,55.
\]
参考整数脚本不得以预存矩阵秩代替上述压缩、关联或割条件。
无 crossing 独立诊断必须输出
\begin{verbatim}
PASS_LOCAL_0_1_COMPRESSION_INEQUALITIES
PASS_CLOSED_NO_CROSSING_TERMINAL_BOUND_AUDIT
shifted_graph_ideals=16
d=9 states=2206 max=35
d=10 states=4057 max=50
d=11 states=7247 max=55
d=12 states=12612 max=60
STATUS=EXACT_INTEGER_DIAGNOSTIC_PASS
\end{verbatim}
闭式边际的独立整数诊断必须输出
\begin{verbatim}
PASS_EXACT_INTEGER_CROSSING_MARGINAL_DENSITY_AUDIT
cut_cost_regimes = 3
active_12_value_marginal_table_required = false
\end{verbatim}

\subsection{shifted 终点的层缺陷规范}

一维 componentwise ideals 恰有十六个。其嵌套最小者的加入次序为
\[
0,1,3,6,2,4,7,5,8,9,
\]
影子剖面为上述 \(n(k)\)。对乘积 ideal 写
\(F_R=\{C:(R,C)\in\mathcal S\}\)、\(k_R=|F_R|\)，并置
\[
m_A=\max_{R\supset A}k_R,\qquad
B(k)=\sum_{A\in\binom{[5]}2}n(m_A).
\]
令 \(\lambda=k^\downarrow\)、
\(H_j=\{R:k_R\ge j\}\)、\(h_j=|H_j|=\lambda'_j\)，并置
\[
e(I)=|\partial I|-n(|I|),\quad
\Delta=(3,2,1,0,2,1,0,1,0,0),\quad
c=(3,2,1,0,2,1,0,1,0,0),
\]
\[
\Phi(\lambda)=\sum_i c_i n(\lambda_i).
\]
审查者应直接复核正文的层缺陷恒等式
\[
B(k)-\Phi(\lambda)=\sum_{j=1}^{10}\Delta_j e(H_j)\ge0.       \tag{LD}
\]
十六 ideals 中仅五个非标准者有非零缺陷；按大小 \(3,4,5,6,7\)
其缺陷依次为 \(1,2,1,1,1\)。由 (LD)：
\begin{enumerate}
\item 同时有至少五行、五列的 Ferrers 分拆，按首行长
  \(5,6,7,8,9,10\) 的最小 \(\Phi\) 依次为
  \(51,54,51,54,54,54\)；
\item 至多四行时，
  \(\Phi(a,b,c,d)=3n(a)+2n(b)+n(c)\)，其不超过五十的解正是
  式~\eqref{n5:eq:fourteen}；连同共轭共十四个；
\item 转置到至多四个实际非空行后，非标准大小四支撑的第一层缺陷
  已为 \(3\cdot2=6\)，故支撑必须为 \(J_4=(0,1,3,6)\)；
\item 对嵌套实际纤维，精确缺陷式为
  \[
  |\partial\mathcal S|-\Phi(\lambda)
  =3e(F_1)+2e(F_2)+e(F_3).
  \]
  余量至多二，按缺陷 \(1,2,1,1,1\) 逐层即强迫全部
  \(F_i=J_{|F_i|}\)。
\end{enumerate}
所以 \(1405\) 个轮廓遍历及七行纤维选择表不再是活跃证明前提。
独立诊断仍应得到
\[
\text{总大小轮廓}=1405,\qquad
\#\{B(k)\le50\}=14,\qquad
\text{非标准低轮廓}=0,
\]
用于交叉核对 (LD)。旧四十九项可见影子表不再是活跃证明前提。

\subsection{五格可见影子的六边父集接口}

对列二元组 \(B\) 重建
\[
P_k(B)=\{c\notin B:B\cup\{c\}\in J_k\}
\]
以及
\[
g_{p,q}(u,v)=\#\{B:P_p(B)\cup P_q(B)\ne\varnothing,\
P_p(B)\subseteq[u],\ P_q(B)\subseteq[v]\}.
\]
对 \(\lambda=(a,b,c,d)\)、面积五分拆
\(\mu=(u_0,\ldots,u_4)\)，独立实现必须逐边得到
\[
\begin{aligned}
h_\lambda(\mu)={}&g_{a,b}(u_2,u_3)+g_{a,c}(u_1,u_3)
 +g_{b,c}(u_1,u_2)\\
&+g_{a,d}(u_0,u_3)+g_{b,d}(u_0,u_2)+g_{c,d}(u_0,u_1).
\end{aligned}
\]
十个 \(B\) 的父列集合必须给出
\[
\begin{array}{c|rrrrrr}
u&0&1&2&3&4&5\\ \hline
g_{3,3}(u,u)&0&3&3&4&6&6\\
g_{4,4}(u,u)&0&0&1&3&6&6\\
g_{5,5}(u,u)&0&1&3&5&7&8\\
g_{10,10}(u,u)&0&0&0&1&4&10.
\end{array}
\]
结合 \(u_1\le2\)、\(u_2,u_3\le1\)，正文式~(VB) 的四个结构上界
必须为 \(7,7,6,6\)，因而相应可见影子下界为
\(41,41,42,42\)。另三族的父集等号必须分别给出：
\[
\begin{array}{c|c|c}
\text{类型}&\max h&\text{等号平面数}\\ \hline
0&20&2\\
1&8&4\\
13&9&2.
\end{array}
\]
这里类型 \(0\) 的两个全局等号平面由坐标转置互换；固定本文取向时
只留下行分拆 \((5)\)。所以 shifted 阈值幸存者仍只有类型
\(0,1,13\)。

独立整数诊断
\path{perm5_visible_shadow_structural_reduction_audit.py}
必须从父点定义重建式~\eqref{n5:eq:six-edge}，不得导入旧四十九项表；
其状态必须为
\texttt{PASS\_EXACT\_INTEGER\_VISIBLE\_SHADOW\_STRUCTURAL\_REDUCTION\_AUDIT}，
且摘要为
\begin{verbatim}
family_profiles                        = 7
plane_partitions                       = 7
diagnostic_pair_values_checked         = 49
five_planes_checked_per_family         = 53130
shifted_survivor_count                 = 3
active_49_value_table_required         = false
non_survivor_uniform_visible_lower_bounds = 41,41,42,42
\end{verbatim}
后两项全遍历只交叉核对六边公式，不承担逻辑责任。数学展开见
\path{n5_visible_shadow_structural_reduction_20260810.md}。

\subsection{类型 1、13 的 Petersen 纤维轮廓接口}

活跃证明不再调用 witness forest。审查者从
\[
(\partial S)_A=N_P\!\left(\bigcup_{R\sim A}S_R\right)
\]
推出
\[
|\partial S|\ge
\sum_A n\!\left(\left|\bigcup_{R\sim A}S_R\right|\right)
\ge
\sum_A n\!\left(\max_{R\sim A}|S_R|\right),
\]
其中
\[
n(k)=(0,3,5,6,6,8,9,9,10,10,10)_k.
\]
一个 ground--set 反射只有三个 Petersen 二轨道，故后一式只含至多
三个整数移动量。审查者不得逐方向抄录旧表；必须先对任意大小轮廓
\(k\) 重建
\[
B(k)-\Phi(k^\downarrow)
=\sum_j\Delta_j\{|N_P(H_j)|-n(|H_j|)\}\ge0,
\qquad H_j=\{R:k_R\ge j\}.                                    \tag{AID}
\]
继而按 Petersen 边次序 \(01,02,03,04,12,13,14,23,24,34\) 重建
\[
\begin{aligned}
r_1&=(0,0,0,2,0,0,4,0,4,10),&
c_1&=(1,1,1,3,1,1,3,1,4,4),\\
r_{13}&=(0,0,0,4,0,0,4,0,5,7),&
c_{13}&=(0,0,1,4,0,1,4,2,4,4).
\end{aligned}
\]
四点邻域最小集必须恰为五个星；部分反射星的邻域必须为八。七点邻域
最小集必须恰为十个 \(N_P(v)^c\)。与十四个低分拆比较后，所有非统一
低 \(B\) 情形必须只剩
\[
\begin{array}{c|c|c}
\text{轮廓与值对}&B\le50\text{ 的计数}&\text{方向}\\ \hline
r_1:(4,2)&0,1,2&r01,r02\\
r_1:(10,2)&0,1,2,3,5,6,7,8&r03\\
c_{13}:(2,0),(1,0)&(0,0),(1,0),(1,1),(2,1)&c03,c13.
\end{array}
\]
这是三种移动模式和五个例外方向；旧四十二个方向记录不是活跃前提。

独立整数诊断
\path{perm5_inverse_shift_layer_orbit_reduction_audit.py}
必须给出
\texttt{PASS\_EXACT\_INTEGER\_INVERSE\_SHIFT\_LAYER\_ORBIT\_REDUCTION\_AUDIT}
及
\begin{verbatim}
base_profiles = 4
diagnostic_unordered_transpositions = 40
stabilizer_transposition_orbit_types = 26
exceptional_movement_schemes = 3
exceptional_transpositions = 5
active_42_direction_table_required = false
remaining_nonuniform_family_states_after_union_bounds = 14
\end{verbatim}
四十个无向换位只交叉核对 (AID) 与三种模式。

五个例外须独立得到：
\[
\begin{array}{c|c|c}
\text{类型、方向}&\text{非统一族状态数}&\text{普通影子大小}\\ \hline
1:r01&2&48\\
1:r02&2&48\\
1:r03&6&48\\
13:c03&2&50\\
13:c13&2&50.
\end{array}
\]
对 \(r03\)，写 \(b=|X\setminus J_4|\)，必须得到
\[
16+2\{n(10-b)+n(4+b)\}
=(48,52,54,52,54,52,48)_b.
\]
对 \(c03,c13\)，必须得到 \(z_0\ne z_2\) 时至少五十二，而
\[
000,010,101,111\longmapsto49,50,50,49.
\]

固定统一族后，按父集
\[
M_q(S)=\{(r,c):(A\cup\{r\},B\cup\{c\})\in S\}
\]
重算 \(h_S(L)=\#\{q:\varnothing\ne M_q(S)\subset L\}\)。五平面极值
必须恰为
\[
\begin{array}{c|rrrrrrr}
&5&4+1&3+2&3+1+1&2+2+1&2+1+1+1&1^5\\ \hline
\max h_{S_1}&7&8&5&5&2&2&0\\
\max h_{S_{13}}&5&6&9&6&7&4&2.
\end{array}
\]
等号平面必须分别只有四个和两个，且各自构成稳定子单轨道。数学展开见
\path{n5_flag_fibre_profile_structural_reduction_20260810.md}。

十四个非统一状态的可见界不得再抄录五行最小值表。置 \(y_i=1-z_i\)，
审查者应直接从 \(M_q(S)\subseteq L\) 重建：
\[
\begin{aligned}
h_{01}&=e_2(z)+z_0z_1z_2+e_2(y)+y_0y_1y_2+y_0+y_1+y_2,\\
h_{02}&=e_3(z_0,z_1,z_2,z_3)+e_3(y_0,y_1,y_2,y_3)
       +z_0y_1+y_0(y_1+y_2+y_3),\\
h_{13}&=7-2z_0-z_1+2z_0z_1 .
\end{aligned}
\]
第一式按 \(|z|=0,1,2,3\) 为 \(7,3,2,4\)；第二式按 \(z_0\) 与
\(z_1+z_2+z_3\) 分开后至多七；第三式按 \(z_0\) 分开后至多七。
对 \(r03\) 的六个状态，必须得到
\[
\begin{gathered}
Q_0=\{12,04;12,14;12,24;12,34\},\\
\mathcal Q_{\{3\}}=Q_0\cup\{23,12;23,13;23,23\},\\
\mathcal Q_{K\cup\{6\}}=Q_0\cup\{02,12;02,13;02,23\},
\end{gathered}
\]
其余四个 \(\mathcal Q_X\) 恰为 \(Q_0\)。故所有状态的被杀父点数
至多七，可见影子分别至少 \(41,43\)。

独立整数诊断
\path{perm5_exceptional_visible_parent_reduction_audit.py}
必须给出
\texttt{PASS\_EXACT\_INTEGER\_EXCEPTIONAL\_VISIBLE\_PARENT\_REDUCTION\_AUDIT}
及
\begin{verbatim}
structural_parent_cases = 3
exceptional_directions = 5
nonuniform_family_states = 14
diagnostic_compatible_family_plane_assignments = 160
maximum_killed_parent_points = 7
uniform_minimum_visible = 41
active_14_state_visible_table_required = false
\end{verbatim}
数学展开见
\path{n5_exceptional_visible_parent_reduction_20260810.md}；160 次代入只
交叉核对上述父集公式。

旧独立整数诊断
\path{perm5_flag_fibre_profile_reduction_audit.py}
不得读取旧森林文件；作为二级交叉核对，其状态仍必须为
\nolinkurl{PASS_EXACT_INTEGER_FIBRE_PROFILE_REDUCTION_AUDIT}，并给出
\begin{verbatim}
directions             = 42
exceptional_directions = 5
nonuniform_states      = 14
minimum_visible        = 41
\end{verbatim}
其中 directions \(=42\) 也只是旧方向层的交叉诊断。旧的
\(123+127=250\) 条 witness 边及 \(3644\) 个局部布尔赋值均不是
证明依赖。数学展开分别见
\path{n5_inverse_shift_layer_orbit_reduction_20260810.md} 与
\path{n5_flag_fibre_profile_structural_reduction_20260810.md}；十四状态的
活跃替代见
\path{n5_exceptional_visible_parent_reduction_20260810.md}。

\subsection{orbit--0 一步逆压缩的 Petersen 接口}

令 \(P=KG(5,2)\)。审查者先证明：若
\(\varnothing\ne X\ne V(P)\)，则
\[
N_P(X)\cap N_P(X^c)\ne\varnothing.
\]
否则长二路端点同色，而 Petersen 图连通且含五圈，迫使所有顶点同色。
这使每个十变量移动块统一；两个同时移动块若统一位相反，则第一坐标
一维影子为六，故乘积影子为六十。

归一化 \(S=S_0\) 后，若 \(L\) 在第零行有 \(5-k\) 格、在第 \(b\) 行
有 \(k\) 格，逐个检查
\[
M_{01}=\{2,3\},\quad
M_{02}=M_{03}=\{1\},\quad
M_{12}=M_{13}=\{0\}
\]
必须给出
\[
|\partial_{L^\perp}S_0|=
\begin{cases}
50-2\binom{k}{3}-2\binom{5-k}{3},&b=1,\\
50-2\binom{5-k}{3},&b=2,3,4.
\end{cases}
\]
由此阈值四十只留下同轨完整行。精确诊断检查 Petersen 图的 1022 个
非平凡子集和全部 31 个有效方向，但不承担证明。

\subsection{orbit--1 长度二接口}

采用正文记号
\[
\begin{aligned}
W_0=\Span\{&
S_{00},S_{01},S_{02},R_{0;01},R_{0;02},R_{0;03},\\
&R_{0;12},R_{0;13},R_{0;23},C_{01;0}\}.
\end{aligned}
\]
在
\[
L=M+\langle x_{10}+Bx_{11}+Cx_{20}+Dx_{21}\rangle
\]
的图表中，固定旗标的非锚顶点恰为二十三个；六个分量连到显式锚，
三个自由三元分量对应 \(B,C,D\)。前九个 \(W_0\) 生成元已经来自
\(q(\Sym^2M)\)。对 \(C_{01;0}\) 的一般提升，依次比较
\[
C_{01;0},X_{01;01},C_{02;0},X_{01;02},X_{01;03},S_{10},
C_{01;1},X_{02;01}
\]
必须得到
\[
a_0=1,\ a_1=B,\ C=0,\ a_2=a_3=a_4=0,\ B^2=0,\ D-BC=0.
\]
反向见证为 \(x_{00}y+Bx_{01}y\)。故完整方案论消元理想为
\[
(B^2,C,D),
\]
而不是只得到相同根集。因此特殊纤维完成局部环为
\(\Q[B]/(B^2)\)，长度二。相对结论使用完成局部环上的拓扑
Nakayama，不得用未经证明的“固定纤维切空间控制移动基”替代。
独立脚本
\path{perm5_orbit1_length2_standalone_exact.py}
只从矩形关系
\(q(x_{ia}x_{jb})=-q(x_{ib}x_{ja})\) 重建上述比较；它没有项目模块
导入，输出
\path{n5_orbit1_length2_standalone_exact.json}，且仅作抄录诊断。

\subsection{orbit--1 十五权纯图终端与 \(W_M\) 接口}

在
\[
L^\dagger=\langle x_{00},x_{01},x_{02},x_{03},x_{10}\rangle
\]
中，商映射在 \(\Sym^2L^\dagger\) 上单射。审查者应直接核查正文
(T1)--(T11)，而不是先运行程序：
\begin{enumerate}
\item 把五个平方记为顶点选择，把十个非平方权记为五顶点图 \(G\)；
\item 对正文的两顶点图 (T5) 和三顶点图 (T6) 逐个删锚、删边，并连同
平方立方块、平方--边块及每个基底三角形的五个逐行换位块，重建
(T4) 的十类、五十个非零 Möbius 项；
\item 把 (T4) 分组为 (T3)，再核查
\(\tau_z+q_4\le6\)、\(\eta\le3\) 与 (T9)；
\item 在 \(H=K_4\) 时以 \(r=c+\rho\) 直接核查 (T10)--(T11) 的
八行，得到最大值三十六且等号恰为 \(W_M,W_1,W_2,W_3\)。
\end{enumerate}

故终端完备性不含计算前提。程序
\path{perm5_orbit1_terminal_pure_formula_audit.py} 只是冗余一致性
诊断：它在全部 \(2^{15}=32768\) 个子集上比较 (T3) 与独立的
\(\Q\) 消元/带符号图引擎，再检查 3003 个十元集。输出必须含
\begin{verbatim}
PASS_EXACT_QQ_REDUNDANT_AUDIT_OF_PURE_GRAPH_FORMULA
all_subsets_formula_vs_QQ_checked = 32768
ten_weight_subsets_checked = 3003
maximum_p = 36
maximum_count = 4
\end{verbatim}
旧脚本 \path{perm5_orbit1_missing_WM_exact.py} 仍应给出
\begin{verbatim}
status = PASS_EXACT_QQ_DIAGNOSTIC_FOUND_MISSING_WM_BRANCH
coordinate_ten_planes_checked = 3003
maximum_p_exact_QQ = 36
p_at_least_36_count = 4
p_WM_exact_QQ = 36
\end{verbatim}
两份输出都只发现抄录错误；任何 PASS 字符串均不替代 (T4) 与
(T7)--(T11) 的正文证明。

\(W_M\) 的局部排除本身是纯证明。作为冗余核对，脚本
\path{perm5_orbit1_WM_same_row_valuative_small_lemmas_QQ_exact.py}
在 \(\Q\) 上检查两个零对角矩阵引理的系数秩为 \(24,25\)，并复算
同行矩形极值序列。它不替代正文的赋值证明。

\subsection{orbit--13 统一结构界接口}

十四权按正文次序编号为 \(0,\ldots,13\)，其中
\(11,12,13\) 是三个 crossing。对固定三次环面权，列出该权的三次
单项式；每个未选择二次商权给一条一元锚或二元带符号边。局部核维数
是未锚且符号相容的连通分量数。

审查者不需要遍历 \(\binom{14}{10}\) 个十元集合。正文结构证明应按
以下四项直接复核：
\begin{enumerate}
\item 对非 crossing 集合写出
\[
p(C)=5\tau+\sum_v(a_v+1)(b_v+1),
\]
并用 \(a_v\le2,b_v\le1\)、端点计费及正文列出的两个剩余小情形证明
\(p(C)\le3|C|\)。
\item 对 crossing \(q=X_{01;ab}\)，选择 \(q\) 按二十五个变量删除
二十五个互异权块中的各一条关系。五个 \(L_A\) 内变量给出总秩降至多
五。
\item 二十个外部变量分成十一项重复行列情形与九项三行三列情形。
前者的两个端点分别由 \(\mathcal U_A\) 外的 square/同行/同列标签
锚定；后者的九条关系是 \(K_{3,3}\)，而所有
\(\mathcal U_A\)-标签边构成一个匹配，删去整个匹配仍为连通六圈。
故外部变量的边际增量为零。
\item 逐次加入至多三个 crossing，得到
\[
p(W)\le3(10-u)+5u\le36<39.
\]
\end{enumerate}

参考实现同时以有理数消元和带符号连通分量核对全部局部 mask，并将
上面的局部形状计数精确复核为 \(5+11+9\)。它还搜索全部
\(2^{14}=16384\) 个子集；该搜索只是反例诊断，不是上述四项证明的
前提。不得调用模三 hypergraph 旧文件。

\subsection{orbit--0 的纯复核接口}

完整旗标的非对角列移动必须由单项方程锚定；二元包含图
\(\binom{[5]}3\leftrightarrow\binom{[5]}2\) 的连通性把未锚问题
降为 \(41\) 个行侧顶点，正文八行表给出八个自由分量。列群在每个
参数上作用平凡，再由完备 Reynolds 与 Nakayama 提升到完整形式局部
环。

Boolean Fourier 复核只需字符正交、射影符号线界和
式~\eqref{n5:eq:shortening}。标准十项余式必须得到
\[
\dim P_3=10,\quad\dim R_3=9,\quad
\dim R_3^{(1)}\le7,\quad
\dim H_3^{(1)}\le35,
\]
从而
\[
\rank\mathcal K_1\ge2215>9\cdot245=2205.
\]
不得把旧的 8008 六元组、Singular 最小素理想或模 101、103 子式
重新列为证明依赖。

\subsection{$n=3,4$ 的活跃核验文件与 SHA--256}

\begingroup\scriptsize
\noindent\path{perm35_exact_verification.py}\\
\nolinkurl{B007564326B8215D730271A30C01D5073024FD3ABEF4CEA394385B8DB021A0E7}\par
\smallskip
\noindent\path{n3_perm35_exact_verification_v11.json}\\
\nolinkurl{87F95D659D1C85EC436D06C66FFAE89C789BAA1A2DBA551417CD1C62A03885F5}\par
\smallskip
\noindent\path{perm4_chow_experiments.py}\\
\nolinkurl{BB6EC23BF6767F589440C44306EECF4A5A94DFF0D526C352B79969A36E2A8B08}\par
\smallskip
\noindent\path{perm4_quadratic_extension_chart_certificate.py}\\
\nolinkurl{4E21096D1A0BFAE70DCB10930886110EB8DA7513A0177C6196347F8D01D04CAB}\par
\smallskip
\noindent\path{perm4_quadratic_extension_chart_exact_verify.py}\\
\nolinkurl{3C1F37913939BA9CF785C494C0F3F15B6A7C5F2A93387A98327D54EF4762DAFB}\par
\smallskip
\noindent\path{perm4_quadratic_extension_full_minor_replay.py}\\
\nolinkurl{BA515FAC7857C2BD4F34B9C061D2B37484266DE6C9D5C7E9BD7C609E92457606}\par
\smallskip
\noindent\path{perm4_rank8_independent_audit.py}\\
\nolinkurl{1EEA3F75D515C2ADDB0550533F7ABA80106EC59D486B3F3A82911EB507C88175}\par
\smallskip
\noindent\path{perm4_rank8_verify_all.py}\\
\nolinkurl{A697830D324F56D9F48002DA0258DD4AF09B9B1456A6BEA0F552918CA03ADA8A}\par
\endgroup

其中 $n=3$ 的 JSON 是程序以素数 $1000003$ 生成的最小秩见证摘要；
$n=4$ 的两个相互独立入口分别重建有理数图表证书，及调用三项紧凑
重放套件。二者刚性结论分别是
\(\rank\mathcal K_1(E_3)=80\)、单项秩 $26$，以及
\(\rank\mathcal K_1(E_4)=560\)、单项秩 $92$、图表常数
\(-32768\)。

\subsection{$n=5$ 的活跃诊断文件与 SHA--256}

\begingroup\scriptsize
\noindent\path{perm5_coordinate_d3_orbit_scan.py}\\
\nolinkurl{6D1B67EE1F99B4BE179C850A57008DEDD086B147237E7C812A838B07F13D82D3}\par
\smallskip
\noindent\path{perm5_coordinate_prolongation_hypergraph.py}\\
\nolinkurl{2AFC8D55CF9BBA4C43D998400F5CDF0F71204974DEDDBB341A2CEC88099EF5CB}\par
\smallskip
\noindent\path{n5_coordinate_prolongation_hypergraph_F3_exact.json}\\
\nolinkurl{584CD155AF8BBB14675AE21EB8E387D809A7C3C7D4977EDCFE6C660B80342DAB}\par
\smallskip
\noindent\path{perm5_fixed_five_shadow_s20_le50_annihilator_classify_exact.py}\\
\nolinkurl{7A101347DEE39995A861D031365AE58C1F0CA9D1B0382615D24CCE95D8F76DAB}\par
\smallskip
\noindent\path{n5_fixed_five_shadow_s20_le50_annihilator_classification_exact.json}\\
\nolinkurl{C06CA7303F31AA1A4BC4125E3CDF03DE066968A70222BCE2F62DC0285D6291A9}\par
\smallskip
\noindent\path{perm5_inverse_shift_layer_orbit_reduction_audit.py}\\
\nolinkurl{31B37F27D95E025296D7CCBED917A71D1B10EAA8C8C5FB675657D810AA8F31D0}\par
\smallskip
\noindent\path{n5_inverse_shift_layer_orbit_reduction_exact.json}\\
\nolinkurl{04A5E08A4F5CBE3621D53C8D0115B159145A0D645626D05262812D75E5DFA3B5}\par
\smallskip
\noindent\path{n5_inverse_shift_layer_orbit_reduction_20260810.md}\\
\nolinkurl{8CF389D88BA5CF657506BF1ABCC1FF234A1586F98C7278DEE2F75ABA454EE95B}\par
\smallskip
\noindent\path{perm5_exceptional_visible_parent_reduction_audit.py}\\
\nolinkurl{FE9098F49AA5ACB02372A330E9CAB02E0504817E7D2C03BE2BB563C9CD515C28}\par
\smallskip
\noindent\path{n5_exceptional_visible_parent_reduction_exact.json}\\
\nolinkurl{F8094A769EB36A888F9971711825FDE541F44C523BB8981273FC49E035D18F4E}\par
\smallskip
\noindent\path{n5_exceptional_visible_parent_reduction_20260810.md}\\
\nolinkurl{E0C8469DA2008692AF432F28322E1DB648A1D5DC0432D1B0042CF2948CC20D02}\par
\smallskip
\noindent\path{perm5_crossing_marginal_density_audit.py}\\
\nolinkurl{5D8EFDFDA7C0CD56C5BA347B09299EE5B7DC05F5B6A84DF9FD573F1AA048454D}\par
\smallskip
\noindent\path{n5_crossing_marginal_density_exact.json}\\
\nolinkurl{F9B0449FAFF8FDD71771C7F48DC81AD65CD03BFFFDB7E2B2BF119974CB966EEB}\par
\smallskip
\noindent\path{n5_crossing_marginal_density_reduction_20260810.md}\\
\nolinkurl{DA0BF22BE942DBDD7FCC6334D587EBF16647EBF53BF8DD2A19ABDB2283304917}\par
\smallskip
\noindent\path{perm5_nocrossing_compression_diagnostic.py}\\
\nolinkurl{8D9EF4E08CEAEB1F3ED9BB315EE71C4124F80AE3BCC5CC12688DCBA242B07AA9}\par
\smallskip
\noindent\path{n5_nocrossing_compression_structural_reduction_20260810.md}\\
\nolinkurl{07C10EFDD807986F463EA06947EEA1CF0F1A90592788FCB629C133F251DFF00E}\par
\smallskip
\noindent\path{perm5_flag_fibre_profile_reduction_audit.py}\\
\nolinkurl{13786E17E9343CFCADFE117178EAC76F26AFE3EE5F3F9A6A5F84406B26F6607E}\par
\smallskip
\noindent\path{n5_flag_fibre_profile_reduction_exact.json}\\
\nolinkurl{5A9B1486F729583F1C77487A9F6B8C8288B387F7EB06DA21AE589069CBC57AD5}\par
\smallskip
\noindent\path{n5_flag_fibre_profile_structural_reduction_20260810.md}\\
\nolinkurl{EA7D3BB63F673881DC6B3E851FAABAFB641FE7DE85D3F3C4DA185F1000927A8E}\par
\smallskip
\noindent\path{perm5_witness_forest_independent_audit.py}\\
\nolinkurl{B20948ED72AFDD13353D8CB68B6015A6BA31B685A528E77C28AEA82CF7F2772F}\par
\smallskip
\noindent\path{n5_witness_forest_independent_audit_exact.json}\\
\nolinkurl{D09B41DEA79243F71039C782772E394C3B0C9A10AD276CCEDD7B5ABDC96F57EF}\par
\smallskip
\noindent\path{perm5_orbit1_length2_standalone_exact.py}\\
\nolinkurl{7759D90B3A4C956D4ACB868FF7FE37DBA1E64A50E3E661923E826F83B20C1EE4}\par
\smallskip
\noindent\path{n5_orbit1_length2_standalone_exact.json}\\
\nolinkurl{8F00D05F30E9566CDA2B500AFC419DE308D89D3475341FA53D3A1EE42C4A4655}\par
\smallskip
\noindent\path{perm5_orbit1_missing_WM_exact.py}\\
\nolinkurl{A6459004F2A53E1FB543F25CE77926A36D5B8402ADF6C72C7E57C90AB0E61E68}\par
\smallskip
\noindent\path{n5_orbit1_missing_WM_exact.json}\\
\nolinkurl{467F84EEFDF0A3D97C50B9A55A6BA6225B957794459EC0808D62F5FDEBD88F07}\par
\smallskip
\noindent\path{n5_orbit1_terminal_pure_graph_classification_20260810.md}\\
\nolinkurl{F79A54CCB8C719888903090CA5DFA22C37539BA2062EF61F2CD108006C9B46FC}\par
\smallskip
\noindent\path{perm5_orbit1_terminal_pure_formula_audit.py}\\
\nolinkurl{A109FD48CD46473BBD82A39A2EAF7501E065052D0D2E02117BAE95C455090A59}\par
\smallskip
\noindent\path{n5_orbit1_terminal_pure_formula_audit_exact.json}\\
\nolinkurl{42FE8F88D51DCC130938125711E99BAEB7F601C107058FFC0DEB079A41B443EA}\par
\smallskip
\noindent\path{perm5_orbit1_WM_same_row_valuative_small_lemmas_QQ_exact.py}\\
\nolinkurl{EF192A6FEC5F66BCB34FE48416C1C34E5DDC10FD808721736C8FE20DDFA02A5E}\par
\smallskip
\noindent\path{n5_orbit1_WM_same_row_valuative_small_lemmas_QQ_exact.json}\\
\nolinkurl{11B4369CFB2072CE3F7BCE8FF4B3E940DF030DB733983D840061BBAC12F770F6}\par
\smallskip
\noindent\path{n5_orbit1_WM_same_row_valuative_two_row_closure_20260810.md}\\
\nolinkurl{1CE28AA9A89194CA93A4465B152DFE6CD095B2E84CA72F7886FE92E6C04B872F}\par
\smallskip
\noindent\path{n5_lower16_route4_orbit1_WM_valuative_splice_20260810.md}\\
\nolinkurl{C4E067FB6A84B3443AD4D6B88014980E3D73904A58613291945D9A7A9315AB18}\par
\smallskip
\noindent\path{n5_lower16_pure_dependency_splice_audit_20260810.md}\\
\nolinkurl{7E8627ACE10D9DB5BEA67C9BC379DFCA26B40332A5F49779C53CCF4E54B6A199}\par
\smallskip
\noindent\path{n5_v13_pure_terminal_self_audit_20260810.md}\\
\nolinkurl{F2B18C968A1B8BFF6CCB5C685EBE7E3A5AA715D0F39E333E2652C67EB47D048F}\par
\smallskip
\noindent\path{perm5_orbit13_four_row_QQ_audit.py}\\
\nolinkurl{5EAD4B3EDD531C4032E0B189183CA70A0FB834DF7FA52431439C10FCEF5DE28F}\par
\smallskip
\noindent\path{n5_orbit13_four_row_QQ_audit_exact.json}\\
\nolinkurl{1703B6BC8C29D2B894FAAA6E2E73CC9272C139AFAB3D54D0507EDD6A9F862DB4}\par
\smallskip
\noindent\path{perm5_orbit13_structural36_audit.py}\\
\nolinkurl{56A5E700C3A4B821CF105D666E6EFB81B69A8703E139F172304C7BDAF7D481B8}\par
\smallskip
\noindent\path{n5_orbit13_structural36_audit_exact.json}\\
\nolinkurl{84E778036B1AE61E3CBB73C5E969DCE813326DC0B7EED5CE3C2507182DAF2566}\par
\smallskip
\noindent\path{n5_orbit13_structural36_pure_20260810.md}\\
\nolinkurl{14ADA341A98E8867248DCA3B5A8B0E30C7647C99273DD39274FD5F5DE799D66E}\par
\smallskip
\noindent\path{perm5_shifted_profile_layer_defect_audit.py}\\
\nolinkurl{3685836E17A8CCC032AE223764F21A3B1D27158BD4A35E7223CA73A0C5E7E10E}\par
\smallskip
\noindent\path{n5_shifted_profile_layer_defect_exact.json}\\
\nolinkurl{D06EE75678C65A4AC32C736150F8E072BD3E9D4721A4173AFD8D233BE003A667}\par
\smallskip
\noindent\path{n5_shifted_profile_layer_defect_reduction_20260810.md}\\
\nolinkurl{B1E6FFCA26774E429EF080C8AC554EBC8309F022AC88F5AC7FA42F9479146B47}\par
\smallskip
\noindent\path{perm5_visible_shadow_structural_reduction_audit.py}\\
\nolinkurl{8E8C5BD638AB30E986F78EB20821B6BE3518AFA9BAE0AF731B97B55FC7974232}\par
\smallskip
\noindent\path{n5_visible_shadow_structural_reduction_exact.json}\\
\nolinkurl{74649A9036A1C4F7BB98C9EB1BA4D10EF37695E8B27B78999B1C5C4ADDF2A23A}\par
\smallskip
\noindent\path{n5_visible_shadow_structural_reduction_20260810.md}\\
\nolinkurl{84BE0315F80FCE40B04545DCF45E5A7A9EB3F5775CF95769B218EA8C45BB520D}\par
\smallskip
\noindent\path{perm5_shifted_terminal_stability_gap_verify.py}\\
\nolinkurl{F17E0E1C32E8BB85646B3F8475A589D05E8DAD7B98B7549B3F870EAF7EE479FB}\par
\smallskip
\noindent\path{n5_shifted_terminal_stability_gap_verify_exact.json}\\
\nolinkurl{123FD0499A399E37C9C4F1265DCD3A0058A75C0DE98C3738D87DF4E58CB85095}\par
\smallskip
\noindent\path{perm5_flag_shifted_stability_verify.py}\\
\nolinkurl{A51B949830BA684970F87F2C41C25206D1DF0453C9884E328096AB8246C03186}\par
\smallskip
\noindent\path{n5_flag_shifted_stability_verify_exact.json}\\
\nolinkurl{CDDFE7D698EC4A0D65369E58761423E926000D1398758D42739461B46DDA3028}\par
\smallskip
\noindent\path{n5_flag_local_witness_forests_20260810.md}\\
\nolinkurl{AB1393F97BE7F4CB40A93CD536EE85D5071670051707673E512D1AA732A1C9EC}\par
\smallskip
\noindent\path{n5_flag_orbit0_full_shadow_witness_forests_20260810.md}\\
\nolinkurl{C59ECF59C54ECCB74101FD39E2AA6A96B50FD4E37A0F6B1A1062EF5E037B86E4}\par
\smallskip
\noindent\path{perm5_orbit0_inverse_shift_petersen_audit.py}\\
\nolinkurl{D14EBE2E15035FE87D179C6630EE0D3076CFDBC4D6DB368564A9845E982D694B}\par
\smallskip
\noindent\path{n5_orbit0_inverse_shift_petersen_exact.json}\\
\nolinkurl{6961C98E312F4949954A8ACDEBBA7E11E5A2600FBD7F38DF247194E99002AB93}\par
\smallskip
\noindent\path{n5_orbit0_inverse_shift_petersen_pure_20260810.md}\\
\nolinkurl{ECEE225790319D732CE580EFF4FB6F56A8B1251658607249737A2D0F1186F226}\par
\endgroup

\subsection{轻量重放顺序}\label{app:replay-order}

\begin{verbatim}
python perm35_exact_verification.py --n 3 --psi
python perm4_rank8_independent_audit.py
python perm4_rank8_verify_all.py
python -O perm5_one_intersection_flag_standalone_exact.py
python perm5_fixed_six_pure_universe_audit.py
python perm5_s19d9_pure_count_audit.py
python perm5_p9_nocrossing_exact.py
python perm5_nocrossing_compression_diagnostic.py
python perm5_crossing_integer_tables_exact.py
python perm5_crossing_marginal_density_audit.py
python perm5_coordinate_prolongation_hypergraph.py
python perm5_p11_global_graph_bound_exact.py
python perm5_p12_global_graph_bound_exact.py
python perm5_shifted_profile_layer_defect_audit.py
python perm5_visible_shadow_structural_reduction_audit.py
python perm5_shifted_terminal_stability_gap_verify.py
python perm5_flag_shifted_stability_verify.py
python perm5_inverse_shift_layer_orbit_reduction_audit.py
python perm5_exceptional_visible_parent_reduction_audit.py
python perm5_flag_fibre_profile_reduction_audit.py
python perm5_orbit0_inverse_shift_petersen_audit.py
python perm5_orbit13_four_row_QQ_audit.py
python perm5_orbit13_structural36_audit.py
python perm5_fixed_five_shadow_s20_le50_annihilator_classify_exact.py
python perm5_s20_orbit1_fixed_flag_P3_graph_exact.py
python perm5_orbit1_length2_standalone_exact.py
python perm5_orbit1_terminal_pure_formula_audit.py
python perm5_orbit1_missing_WM_exact.py
python perm5_orbit1_WM_same_row_valuative_small_lemmas_QQ_exact.py
python perm5_s20_orbit0_fullflag_tangent_graph_exact.py
python perm5_orbit0_fourier_rigidity_verify.py
\end{verbatim}
其中 \path{perm5_coordinate_d3_orbit_scan.py} 仅作为 $p_{11},p_{12}$
及局部表生成器导入的确定性坐标类构造模块；轻量重放不执行其历史
有限域扫描入口。局部表生成器的 $\mathbb F_3$ 输出仅核对正文纯整数割
引理的两个极端值，不进入定理的逻辑依赖。
每个生产者先覆盖自己的声明输出，再立即检查状态、计数和 SHA--256。
预期总状态为
\begin{verbatim}
PASS_STRUCTURAL_REDUCTION_LOWER16_INTERNAL_AUDIT
shifted_profiles          = 14
active_1405_profile_enum  = 0
active_fibre_choice_table = 0
active_49_visible_table   = 0
visible_structural_bounds = 4
active_42_direction_table = 0
exceptional_move_schemes  = 3
family_exception_dirs     = 5
nonuniform_family_states  = 14
active_14_state_table     = 0
exceptional_parent_cases  = 3
active_nocrossing_square_orbit_tables = 0
nocrossing_closed_exception_types = 10
active_12_crossing_table  = 0
crossing_density_regimes  = 3
global_good_planes        = 6
active_witness_edges      = 0
orbit0_witness_edges      = 0
orbit13_structural_bound = 36
orbit13_subsets_required = 0
orbit1_terminal_subsets_required = 0
orbit1_terminal_mobius_families = 10
orbit1_terminal_nonzero_terms = 50
orbit1_terminal_K4_rows = 8
orbit1_terminal_max_p = 36
orbit1_terminal_max_count = 4
orbit1_terminal_QQ_role = diagnostic
one_intersection_flag_role = theorem_premise
one_intersection_flags = 886464
one_intersection_max_F3 = 26
old_10GB_assets_needed = false
\end{verbatim}
PDF 还内嵌本 LaTeX 源、层缺陷笔记、六边父集笔记、十四状态父集公式、
crossing 边际密度证明、Petersen 纤维结构笔记、两份旧 witness
交叉诊断附表、一维交旗标独立重建程序及其紧凑 JSON、$n=3,4,5$ 的
全部上述核验程序以及相应的最小 JSON。
历史 SAT/DRAT、60,960 旗标闭包、1096 标准集、有限域矩阵、Singular
输出和 10GB 文件均不应出现在重放拓扑或 PDF 附件中。
